% ------------------------------------------------------------------------------
% Michael Kastl 87 28 60
% ------------------------------------------------------------------------------
%
%
% Dokumentenkopf 
% ------------------------------------------------------------------------------
\documentclass[
    11pt, % Schriftgr��e
    DIV10,
    ngerman, % f�r Umlaute, Silbentrennung etc.
    a4paper, % Papierformat
    oneside, % einseitiges Dokument
    titlepage, % es wird eine Titelseite verwendet
    parskip=half, % Abstand zwischen Abs�tzen (halbe Zeile)
    headings=normal, % Gr��e der �berschriften verkleinern
    listof=totoc, % Verzeichnisse im Inhaltsverzeichnis auff�hren
    bibliography=totoc, % Literaturverzeichnis im Inhaltsverzeichnis auff�hren
    index=totoc, % Index im Inhaltsverzeichnis auff�hren
    captions=tableheading, % Beschriftung von Tabellen unterhalb ausgeben
    final, % Status des Dokuments (final/draft)
    numbers=noenddot
]{scrreprt}

\renewcommand*{\chapterheadstartvskip}{\vspace{-1\baselineskip}} %Abstand vor Section
\setlength{\headsep}{30pt} %Abstand Kopfzeile Textk�rper


%Abstand vor und nach �berschrift einstellbar
%\makeatletter
%\def\@makeschapterhead#1{%
% \vspace*{5\p@}%       % 5 = Abstand vor chapter
%  {\parindent \z@ \raggedright
%    \normalfont
%    \interlinepenalty\@M
%    \huge \bfseries  #1\par\nobreak 
%    \vskip 30\p@		% 30 = Abstand nach chapter
%  }}
%\makeatother


% Meta-Informationen -----------------------------------------------------------
%   Informationen �ber das Dokument, wie z.B. Titel, Autor, Matrikelnr. etc
%   werden in der Datei Meta.tex definiert und k\"onnen danach global
%   verwendet werden.
% ------------------------------------------------------------------------------
% Meta-Informationen -----------------------------------------------------------
%   Definition von globalen Parametern, die im gesamten Dokument verwendet
%   werden k�nnen (z.B auf dem Deckblatt etc.).
%
%   ACHTUNG: Wenn die Texte Umlaute oder ein Esszet enthalten, muss der folgende
%            Befehl bereits an dieser Stelle aktiviert werden:
%            \usepackage[latin1]{inputenc}
% ------------------------------------------------------------------------------
\newcommand{\titel}{Vorlage f\"ur eine Diplomarbeit in LaTeX}
\newcommand{\untertitel}{\ }
\newcommand{\art}{Diplomarbeit}
\newcommand{\fachgebiet}{Automatisierungstechnik}
\newcommand{\autor}{Max Mustermann}
\newcommand{\studienbereich}{Mechatronik}
\newcommand{\matrikelnr}{xx xx xx}
\newcommand{\erstgutachter}{Prof. Dr.-Ing. Max Mustermann}
\newcommand{\zweitgutachter}{Dipl.-Ing. xxx}
\newcommand{\jahr}{2010}
\newcommand{\ort}{M"unchen}
\newcommand{\logo}{mkDoc-Logo.png}


% ben\"otigte Packages -----------------------------------------------------------
%   LaTeX-Packages, die ben�tigt werden, sind in die Datei Packages.tex
%   "ausgelagert", um diese Vorlage m�glichst �bersichtlich zu halten.
% ------------------------------------------------------------------------------
% Anpassung des Seitenlayouts --------------------------------------------------
%   siehe Seitenstil.tex
% ------------------------------------------------------------------------------
\usepackage[
    automark, % Kapitelangaben in Kopfzeile automatisch erstellen
    headsepline, % Trennlinie unter Kopfzeile
    ilines % Trennlinie linksb�ndig ausrichten
]{scrpage2}

% Anpassung an Landessprache ---------------------------------------------------
\usepackage[ngerman]{babel}

% Umlaute ----------------------------------------------------------------------
%   Umlaute/Sonderzeichen wie ���� direkt im Quelltext verwenden (CodePage).
%   Erlaubt automatische Trennung von Worten mit Umlauten.
% ------------------------------------------------------------------------------
\usepackage[latin1]{inputenc}
\usepackage[T1]{fontenc}
\usepackage{textcomp} % Euro-Zeichen etc.

% Schrift ----------------------------------------------------------------------
\usepackage{lmodern} % bessere Fonts
\usepackage{relsize} % Schriftgr��e relativ festlegen
%\usepackage{subscript}
\newcommand{\mysub}[1]{\raisebox{-0.5ex}{\scriptsize{#1}}}

%Formeln
\usepackage{nicefrac}

% Grafiken ---------------------------------------------------------------------
% Einbinden von JPG-Grafiken erm�glichen
\usepackage[dvips,final]{graphicx}
% hier liegen die Bilder des Dokuments
\graphicspath{{Bilder/}}

% Befehle aus AMSTeX f�r mathematische Symbole z.B. \boldsymbol \mathbb --------
\usepackage{amsmath,amsfonts}

% f�r Index-Ausgabe mit \printindex --------------------------------------------
\usepackage{makeidx}

% Einfache Definition der Zeilenabst�nde und Seitenr�nder etc. -----------------
\usepackage{setspace}
\usepackage{geometry}

% Symbolverzeichnis ------------------------------------------------------------
%   Symbolverzeichnisse bequem erstellen. Beruht auf MakeIndex:
%     makeindex.exe %Name%.nlo -s nomencl.ist -o %Name%.nls
%   erzeugt dann das Verzeichnis. Dieser Befehl kann z.B. im TeXnicCenter
%   als Postprozessor eingetragen werden, damit er nicht st�ndig manuell
%   ausgef�hrt werden muss.
%   Die Definitionen sind ausgegliedert in die Datei "Glossar.tex".
% ------------------------------------------------------------------------------
\usepackage[printonlyused]{acronym}

% zum Umflie�en von Bildern ----------------------------------------------------
%\usepackage{floatflt}


% zum Einbinden von Programmcode -----------------------------------------------
\usepackage{listings}
\usepackage{xcolor} 
\definecolor{hellgelb}{rgb}{1,1,0.9}
\definecolor{colKeys}{rgb}{0,0,1}
\definecolor{colIdentifier}{rgb}{0,0,0}
\definecolor{colComments}{rgb}{1,0,0}
\definecolor{colString}{rgb}{0,0.5,0}
\lstset{
    float=hbp,
    basicstyle=\ttfamily\color{black}\small\smaller,
    identifierstyle=\color{colIdentifier},
    keywordstyle=\color{colKeys},
    stringstyle=\color{colString},
    commentstyle=\color{colComments},
    columns=flexible,
    tabsize=2,
    frame=single,
    extendedchars=true,
    showspaces=false,
    showstringspaces=false,
    numbers=left,
    numberstyle=\tiny,
    breaklines=true,
    backgroundcolor=\color{hellgelb},
    breakautoindent=true
}

% URL verlinken, lange URLs umbrechen etc. -------------------------------------
\usepackage{url}

% wichtig f�r korrekte Zitierweise ---------------------------------------------
\usepackage[square]{natbib}

%Farben definieren
\definecolor{darkblue}{rgb}{0,0,.5}

% PDF-Optionen -----------------------------------------------------------------

\usepackage[
    bookmarks,
    bookmarksopen=true,
    colorlinks=true,
% diese Farbdefinitionen zeichnen Links im PDF farblich aus
    linkcolor=blue, % einfache interne Verkn�pfungen
    anchorcolor=black,% Ankertext
    citecolor=blue, % Verweise auf Literaturverzeichniseintr�ge im Text
    filecolor=magenta, % Verkn�pfungen, die lokale Dateien �ffnen
    menucolor=red, % Acrobat-Men�punkte
    urlcolor=cyan, 
% diese Farbdefinitionen sollten f�r den Druck verwendet werden (alles schwarz)
    %linkcolor=black, % einfache interne Verkn�pfungen
    %anchorcolor=black, % Ankertext
    %citecolor=black, % Verweise auf Literaturverzeichniseintr�ge im Text
    %filecolor=black, % Verkn�pfungen, die lokale Dateien �ffnen
    %menucolor=black, % Acrobat-Men�punkte
    %urlcolor=black, 
    backref,
    plainpages=false, % zur korrekten Erstellung der Bookmarks
    pdfpagelabels, % zur korrekten Erstellung der Bookmarks
    hypertexnames=false, % zur korrekten Erstellung der Bookmarks
    linktocpage % Seitenzahlen anstatt Text im Inhaltsverzeichnis verlinken
]{hyperref}
% Befehle, die Umlaute ausgeben, f�hren zu Fehlern, wenn sie hyperref als Optionen �bergeben werden
\hypersetup{
    pdftitle={\titel \untertitel},
    pdfauthor={\autor},
    pdfcreator={\autor},
    pdfsubject={\titel \untertitel},
    pdfkeywords={\titel \untertitel},
}

% fortlaufendes Durchnummerieren der Fu�noten ----------------------------------
\usepackage{chngcntr}

% f�r lange Tabellen -----------------------------------------------------------
\usepackage{longtable}
\usepackage{colortbl}
\usepackage{array}
\usepackage{ragged2e}
\usepackage{lscape}

% Spaltendefinition rechtsb�ndig mit definierter Breite ------------------------
\newcolumntype{w}[1]{>{\raggedleft\hspace{0pt}}p{#1}}

% Formatierung von Listen �ndern -----------------------------------------------
\usepackage{paralist}

% bei der Definition eigener Befehle ben�tigt
\usepackage{ifthen}

% definiert u.a. die Befehle \todo und \listoftodos
\usepackage{todonotes}

% sorgt daf�r, dass Leerzeichen hinter parameterlosen Makros nicht als Makroendezeichen interpretiert werden
\usepackage{xspace}


% Erstellung eines Index und Abk�rzungsverzeichnisses aktivieren ---------------

\makeindex
%\makenomenclature

% Kopf- und Fu�zeilen, Seitenr�nder etc. ---------------------------------------
% Zeilenabstand 1,5 Zeilen -----------------------------------------------------
\onehalfspacing

% Seitenr�nder -----------------------------------------------------------------
\setlength{\topskip}{\ht\strutbox} % behebt Warnung von geometry
\geometry{paper=a4paper,left=35mm,right=35mm,top=30mm}

% Kopf- und Fu�zeilen ----------------------------------------------------------
\pagestyle{scrheadings}
% Kopf- und Fu�zeile auch auf Kapitelanfangsseiten
\renewcommand*{\chapterpagestyle}{scrheadings} 
% Schriftform der Kopfzeile
\renewcommand{\headfont}{\normalfont}

% Kopfzeile
\ihead{\large{\textsc{\titel}}\\[2ex] \textit{\headmark}}
\chead{}
\ohead{\includegraphics[scale=0.3]{\logo}}
\setlength{\headheight}{21mm} % H�he der Kopfzeile
% Kopfzeile �ber den Text hinaus verbreitern
\setheadwidth[0pt]{textwithmarginpar} 
\setheadsepline[text]{0.4pt} % Trennlinie unter Kopfzeile

% Fu�zeile
\ifoot{\copyright\ \autor}
\cfoot{}
\ofoot{\pagemark}

% sonstige typographische Einstellungen ----------------------------------------

% erzeugt ein wenig mehr Platz hinter einem Punkt
\frenchspacing 

% Schusterjungen und Hurenkinder vermeiden
\clubpenalty = 10000
\widowpenalty = 10000 
\displaywidowpenalty = 10000

% Quellcode-Ausgabe formatieren
\lstset{numbers=left, numberstyle=\tiny, numbersep=5pt, breaklines=true}
\lstset{emph={square}, emphstyle=\color{red}, emph={[2]root,base}, emphstyle={[2]\color{blue}}}

% Fu�noten fortlaufend durchnummerieren
\counterwithout{footnote}{chapter}


% eigene Definitionen f�r Silbentrennung
% Trennvorschl�ge im Text werden mit \" angegeben
% untrennbare W�rter und Ausnahmen von der normalen Trennung k�nnen in dieser
% Datei mittels \hyphenation definiert werden


% eigene LaTeX-Befehle
% Eigene Befehle und typographische Auszeichnungen f�r diese

% einfaches Wechseln der Schrift, z.B.: \changefont{cmss}{sbc}{n}
\newcommand{\changefont}[3]{\fontfamily{#1} \fontseries{#2} \fontshape{#3} \selectfont}

% Abk�rzungen mit korrektem Leerraum 
\newcommand{\ua}{\mbox{u.\,a.\ }}
\newcommand{\zB}{\mbox{z.\,B.\ }}
\newcommand{\dahe}{\mbox{d.\,h.\ }}
\newcommand{\Vgl}{Vgl.\ }
\newcommand{\bzw}{bzw.\ }
\newcommand{\evtl}{evtl.\ }

\newcommand{\abbildung}[1]{Abbildung~\ref{fig:#1}}

\newcommand{\bs}{$\backslash$}

% erzeugt ein Listenelement mit fetter �berschrift 
\newcommand{\itemd}[2]{\item{\textbf{#1}}\\{#2}}

% einige Befehle zum Zitieren --------------------------------------------------
\newcommand{\Zitat}[2][\empty]{\ifthenelse{\equal{#1}{\empty}}{\citep{#2}}{\citep[#1]{#2}}}

% zum Ausgeben von Autoren
\newcommand{\AutorName}[1]{\textsc{#1}}
\newcommand{\Autor}[1]{\AutorName{\citeauthor{#1}}}

% verschiedene Befehle um W�rter semantisch auszuzeichnen ----------------------
\newcommand{\NeuerBegriff}[1]{\textbf{#1}}
\newcommand{\Fachbegriff}[1]{\textit{#1}}

\newcommand{\Eingabe}[1]{\texttt{#1}}
\newcommand{\Code}[1]{\texttt{#1}}
\newcommand{\Datei}[1]{\texttt{#1}}

\newcommand{\Datentyp}[1]{\textsf{#1}}
\newcommand{\XMLElement}[1]{\textsf{#1}}
\newcommand{\Webservice}[1]{\textsf{#1}}


% Das eigentliche Dokument -----------------------------------------------------
%   Der eigentliche Inhalt des Dokuments beginnt hier. Die einzelnen Seiten
%   und Kapitel werden in eigene Dateien ausgelagert und hier nur inkludiert.
% ------------------------------------------------------------------------------
\begin{document}

% auch subsubsection nummerieren
\setcounter{secnumdepth}{3}
\setcounter{tocdepth}{3}

% Deckblatt und Abstract ohne Seitenzahl
\ofoot{}
\thispagestyle{plain}
\begin{titlepage}

\begin{center}

\huge{\textbf{\titel}}\\[1.5ex]
\LARGE{\textbf{\untertitel}}\\[6ex]
\LARGE{\textbf{\art}}\\[1.5ex]
\Large{im Fachgebiet \fachgebiet}\\[12ex]

\includegraphics[scale=1]{mkDoc-Logo.png}\\[14ex]

\normalsize
\begin{tabular}{w{5.4cm}p{6cm}}\\
vorgelegt von:  & \quad \autor\\[1.2ex]
Studienbereich: & \quad \studienbereich\\[1.2ex]
Matrikelnummer: & \quad \matrikelnr\\[1.2ex]
Erstgutachter:  & \quad \erstgutachter\\[4ex]
\end{tabular}

\copyright\ \jahr\\[9ex]

\end{center}

%\singlespacing
%\small
%\noindent Dieses Werk einschlie�lich seiner Teile ist \textbf{urheberrechtlich gesch�tzt}. Jede Verwertung au�erhalb der engen Grenzen des Urheberrechtgesetzes ist ohne Zustimmung des Autors unzul�ssig und strafbar. Das gilt insbesondere f�r Vervielf�ltigungen, �bersetzungen, Mikroverfilmungen sowie die Einspeicherung und Verarbeitung in elektronischen Systemen.

\end{titlepage}

\include{Inhalt/Abstract}
\ofoot{\pagemark}

% Seitennummerierung -----------------------------------------------------------
%   Vor dem Hauptteil werden die Seiten in gro�en r�mischen Ziffern 
%   nummeriert.
% ------------------------------------------------------------------------------
\pagenumbering{Roman}
\tableofcontents % Inhaltsverzeichnis

% Abk�rzungsverzeichnis --------------------------------------------------------

%\chapter*
\addchap{Abk\"urzungsverzeichnis}
\begin{acronym}[laaaaaaaaaaang] %In der Eckigen Klammer muss das l�ngste Akronym stehen. Es wird dann die gesamte weitere Tabelle danach ausgerichtet
  \acro{AM}{andere Ma\ss nahmen}
  \acro{B10d}[$\text{B}_{\text{10d}}$]{Anzahl von Zyklen, bis 10 \% der Komponenten gef\"ahrlich ausgefallen sind (f\"ur pneumatische und elektromechanische Komponenten)}
  \acro{CCF}{Ausfall aufgrund gemeinsamer Ursache}
  \acro{DC}{Diagnosedeckungsgrad}
  \acro{DCavg}[$\text{DC}_{\text{avg}}$]{durchschnittlicher Diagnosedeckungsgrad}
  \acro{dop}[$\text{d}_{\text{op}}$]{mittlere Betriebszeit in Tagen je Jahr}
  \acro{EMC}{Elektromagnetische Vertr\"aglichkeit}
  \acro{FMEA}{Fehlzustandsart- und -auswirkungsanalyse}
  \acro{hop}[$\text{h}_{\text{op}}$]{mittlere Betriebszeit in Stunden je Tag}
  \acro{MTTF}{mittlere Zeit bis zum Ausfall}
  \acro{MTTFd}[$\text{MTTF}_{\text{d}}$]{mittlere Zeit bis zum gefahrbringenden Ausfall}
  \acro{MTTFdC1}[$\text{MTTF}_{\text{dC1}}$]{mittlere Zeit bis zum gefahrbringenden Ausfall Kanal 1}
  \acro{MTTFdC2}[$\text{MTTF}_{\text{dC2}}$]{mittlere Zeit bis zum gefahrbringenden Ausfall Kanal 2}
  \acro{MTTFdi}[$\text{MTTF}_{\text{di}}$]{mittlere Zeit bis zum gefahrbringenden Ausfall eines Bauteils}
  \acro{nop}[$\text{n}_{\text{op}}$]{mittlere Anzahl j\"ahrlicher Bet\"atigungen}
  \acro{PFHD}[$\text{PFH}_{\text{d}}$]{Wahrscheinlichkeit eines gefahrbringenden Ausfalls pro Stunde}
  \acro{PL}{Performance Level}
  \acro{PLr}[$\text{PL}_{\text{r}}$]{erforderliches Performance Level}
  \acro{SIL}{Safety integrated Level}
  \acro{SLS}{sicher begrenzte Geschwindigkeit}
  \acro{SOS}{Stillstands\"uberwachung}
  \acro{SRCF}{sicherheitsbezogene Steuerungsfunktion}
  \acro{SRECS}{sicherheitsbezogenes elektrisches Steuerungssystem}
  \acro{SRP/CS}{sicherheitsbezogenes Teil einer Steuerung}
  \acro{STO}{sicherer Halt} 
  \acro{tzyklus}[$\text{t}_{\text{Zyklus}}$]{mittlere Zeit zwischen dem Beginn zweier aufeinander folgenden Zyklen des Bauteils (z. B. Schalten eines Ventils) in Sekunden je Zyklus}
  \acro{VDMA}{Verband Deutscher Maschinen- und Anlagenbau}
  \acro{lambdad}[$\lambda_{\text{d}}$]{Anzahl aller gef\"ahrlichen Ausf\"alle}
  \acro{lambdadd}[$\lambda_{\text{dd}}$]{Anzahl aller erkannter gef\"ahrlichen Ausf\"alle}
\end{acronym}

\listoffigures % Abbildungsverzeichnis
\listoftables % Tabellenverzeichnis
%\renewcommand{\lstlistlistingname}{Verzeichnis der Listings}
%\lstlistoflistings % Listings-Verzeichnis

% arabische Seitenzahlen im Hauptteil ------------------------------------------
\clearpage
\pagenumbering{arabic}

% die Inhaltskapitel werden in "Inhalt.tex" inkludiert -------------------------
% Hier k�nnen die einzelnen Kapitel inkludiert werden. Sie m�ssen in den 
% entsprechenden .TEX-Dateien vorliegen. Die Dateinamen k�nnen nat�rlich 
% angepasst werden.
%%\chapter*
\addchap{Abk\"urzungsverzeichnis}
\begin{acronym}[laaaaaaaaaaang] %In der Eckigen Klammer muss das l�ngste Akronym stehen. Es wird dann die gesamte weitere Tabelle danach ausgerichtet
  \acro{AM}{andere Ma\ss nahmen}
  \acro{B10d}[$\text{B}_{\text{10d}}$]{Anzahl von Zyklen, bis 10 \% der Komponenten gef\"ahrlich ausgefallen sind (f\"ur pneumatische und elektromechanische Komponenten)}
  \acro{CCF}{Ausfall aufgrund gemeinsamer Ursache}
  \acro{DC}{Diagnosedeckungsgrad}
  \acro{DCavg}[$\text{DC}_{\text{avg}}$]{durchschnittlicher Diagnosedeckungsgrad}
  \acro{dop}[$\text{d}_{\text{op}}$]{mittlere Betriebszeit in Tagen je Jahr}
  \acro{EMC}{Elektromagnetische Vertr\"aglichkeit}
  \acro{FMEA}{Fehlzustandsart- und -auswirkungsanalyse}
  \acro{hop}[$\text{h}_{\text{op}}$]{mittlere Betriebszeit in Stunden je Tag}
  \acro{MTTF}{mittlere Zeit bis zum Ausfall}
  \acro{MTTFd}[$\text{MTTF}_{\text{d}}$]{mittlere Zeit bis zum gefahrbringenden Ausfall}
  \acro{MTTFdC1}[$\text{MTTF}_{\text{dC1}}$]{mittlere Zeit bis zum gefahrbringenden Ausfall Kanal 1}
  \acro{MTTFdC2}[$\text{MTTF}_{\text{dC2}}$]{mittlere Zeit bis zum gefahrbringenden Ausfall Kanal 2}
  \acro{MTTFdi}[$\text{MTTF}_{\text{di}}$]{mittlere Zeit bis zum gefahrbringenden Ausfall eines Bauteils}
  \acro{nop}[$\text{n}_{\text{op}}$]{mittlere Anzahl j\"ahrlicher Bet\"atigungen}
  \acro{PFHD}[$\text{PFH}_{\text{d}}$]{Wahrscheinlichkeit eines gefahrbringenden Ausfalls pro Stunde}
  \acro{PL}{Performance Level}
  \acro{PLr}[$\text{PL}_{\text{r}}$]{erforderliches Performance Level}
  \acro{SIL}{Safety integrated Level}
  \acro{SLS}{sicher begrenzte Geschwindigkeit}
  \acro{SOS}{Stillstands\"uberwachung}
  \acro{SRCF}{sicherheitsbezogene Steuerungsfunktion}
  \acro{SRECS}{sicherheitsbezogenes elektrisches Steuerungssystem}
  \acro{SRP/CS}{sicherheitsbezogenes Teil einer Steuerung}
  \acro{STO}{sicherer Halt} 
  \acro{tzyklus}[$\text{t}_{\text{Zyklus}}$]{mittlere Zeit zwischen dem Beginn zweier aufeinander folgenden Zyklen des Bauteils (z. B. Schalten eines Ventils) in Sekunden je Zyklus}
  \acro{VDMA}{Verband Deutscher Maschinen- und Anlagenbau}
  \acro{lambdad}[$\lambda_{\text{d}}$]{Anzahl aller gef\"ahrlichen Ausf\"alle}
  \acro{lambdadd}[$\lambda_{\text{dd}}$]{Anzahl aller erkannter gef\"ahrlichen Ausf\"alle}
\end{acronym}

\chapter{Einleitung}
\label{cha:Einleitung}

\section{Section 1}
\label{sec:1Section1}
Willkommen im Portal f"ur Elektronik, Maschinenbau und Mechatronik !
Dieses Portal soll euch beim lernen und diskutieren der einzelnen Studienf"acher behilflich sein, oder im Alltag als Knowledge-Base zur Verf"ugung stehen ! \\
F"ur jedes Studienfach wird in einem "Ubersichtsartikel der Inhalt zusammengefasst und die einzelnen Fachartikel in Beziehung zueinader gestellt. Kommen Formeln in den Fachartikeln vor, werden diese in einer Formelsammlung zu dem jeweiligen Studienfach zusammengef"uhrt. Am Ende soll jedes Studienfach einen "Ubersichtsartikel und wenn m"oglich eine Formelsammlung besitzen.

\subsection{Subcestion 1.1}
\label{subsec:1Subcestion1.1}

\begin{figure}[htb]
\centering
\includegraphics[width=0.2\textwidth]{mkDoc-Logo.png}
\caption{mkDoc}
\label{fig:mkDoc}
\end{figure}


\subsection{Subcestion 1.2}
\label{subsec:1Subcestion 1.2}
Willkommen im Portal f"ur Elektronik, Maschinenbau und Mechatronik !\footnote{\Vgl\Zitat[S.~11]{Sicherheitstechnik}}
Dieses Portal soll euch beim lernen und diskutieren der einzelnen Studienf"acher behilflich sein, oder im Alltag als Knowledge-Base zur Verf"ugung stehen ! \\
F"ur jedes Studienfach wird in einem "Ubersichtsartikel der Inhalt zusammengefasst und die einzelnen Fachartikel in Beziehung zueinader gestellt. Kommen Formeln in den Fachartikeln vor, werden diese in einer Formelsammlung zu dem jeweiligen Studienfach zusammengef"uhrt. Am Ende soll jedes Studienfach einen "Ubersichtsartikel und wenn m"oglich eine Formelsammlung besitzen.

\chapter{Chapter 2}
\label{cha:Chapter2}

\section{Section 1}
\label{sec:2Section1}
Willkommen im Portal f"ur Elektronik, Maschinenbau und Mechatronik !
Dieses Portal soll euch beim lernen und diskutieren der einzelnen Studienf"acher behilflich sein, oder im Alltag als Knowledge-Base zur Verf"ugung stehen ! \\
F"ur jedes Studienfach wird in einem "Ubersichtsartikel der Inhalt zusammengefasst und die einzelnen Fachartikel in Beziehung zueinader gestellt. Kommen Formeln in den Fachartikeln vor, werden diese in einer Formelsammlung zu dem jeweiligen Studienfach zusammengef"uhrt. Am Ende soll jedes Studienfach einen "Ubersichtsartikel und wenn m"oglich eine Formelsammlung besitzen.

\subsection{Subcestion 1.1}
\label{subsec:2Subcestion1.1}

\begin{figure}[htb]
\centering
\includegraphics[width=0.2\textwidth]{mkDoc-Logo.png}
\caption{mkDoc}
\label{fig:mkDoc}
\end{figure}


\subsection{Subcestion 1.2}
\label{subsec:2Subcestion 1.2}
Willkommen im Portal f"ur Elektronik, Maschinenbau und Mechatronik !\footnote{\Vgl\Zitat[S.~11]{Sicherheitstechnik}}
Dieses Portal soll euch beim lernen und diskutieren der einzelnen Studienf"acher behilflich sein, oder im Alltag als Knowledge-Base zur Verf"ugung stehen ! \\
F"ur jedes Studienfach wird in einem "Ubersichtsartikel der Inhalt zusammengefasst und die einzelnen Fachartikel in Beziehung zueinader gestellt. Kommen Formeln in den Fachartikeln vor, werden diese in einer Formelsammlung zu dem jeweiligen Studienfach zusammengef"uhrt. Am Ende soll jedes Studienfach einen "Ubersichtsartikel und wenn m"oglich eine Formelsammlung besitzen.

\section{Section 2}
\label{sec:2Section2}

Willkommen im Portal f"ur Elektronik, Maschinenbau und Mechatronik !
Dieses Portal soll euch beim lernen und diskutieren der einzelnen Studienf"acher behilflich sein, oder im Alltag als Knowledge-Base zur Verf"ugung stehen ! \ref{sec:2Section1} \\
F"ur jedes Studienfach wird in einem "Ubersichtsartikel der Inhalt zusammengefasst und die einzelnen Fachartikel in Beziehung zueinader gestellt. Kommen Formeln in den Fachartikeln vor, werden diese in einer Formelsammlung zu dem jeweiligen Studienfach zusammengef"uhrt. Am Ende soll jedes Studienfach einen "Ubersichtsartikel und wenn m"oglich eine Formelsammlung besitzen.

\section{Section 3}
\label{sec:2Section3}

\subsection{Subcestion 3.1}
\label{subsec:2Subcestion3.1}
Willkommen im Portal f"ur Elektronik, Maschinenbau und Mechatronik !
Dieses Portal soll euch beim lernen und diskutieren der einzelnen Studienf"acher behilflich sein, oder im Alltag als Knowledge-Base zur Verf"ugung stehen ! \\
F"ur jedes Studienfach wird in einem "Ubersichtsartikel der Inhalt zusammengefasst und die einzelnen Fachartikel in Beziehung zueinader gestellt. Kommen Formeln in den Fachartikeln vor, werden diese in einer Formelsammlung zu dem jeweiligen Studienfach zusammengef"uhrt. Am Ende soll jedes Studienfach einen "Ubersichtsartikel und wenn m"oglich eine Formelsammlung besitzen. \acs{MTTF}

\subsection{Subcestion 3.2}
\label{subsec:2Subcestion3.2}
Willkommen im Portal f"ur Elektronik, Maschinenbau und Mechatronik !
Dieses Portal soll euch beim lernen und diskutieren der einzelnen Studienf"acher behilflich sein, oder im Alltag als Knowledge-Base zur Verf"ugung stehen ! \\
F"ur jedes Studienfach wird in einem "Ubersichtsartikel der Inhalt zusammengefasst und die einzelnen Fachartikel in Beziehung zueinader gestellt. Kommen Formeln in den Fachartikeln vor, werden diese in einer Formelsammlung zu dem jeweiligen Studienfach zusammengef"uhrt. Am Ende soll jedes Studienfach einen "Ubersichtsartikel und wenn m"oglich eine Formelsammlung besitzen.

\newpage
\subsection{Subcestion 3.3}
\label{subsec:2Subcestion3.3}

\begin{longtable}{|p{9cm}|c|}
\caption{Klassifikation der Schwere (S)}\footnote{\Vgl\Zitat[S.~76]{DINEN62061}}\\
\hline
\label{tab:KlassifikationSchwere}
\textbf{Auswirkungen} & \textbf{Schwere (S)}\\
\hline
\hline
irreversibel: Tod, Verlust eines Auges oder Arms & 4\\
\hline
irreversibel: gebrochene Gliedma\ss en, Verlust (eines) mehrerer Finger(s) & 3\\
\hline
reversibel: Behandlung durch einen Mediziner erforderlich & 2\\
\hline
reversibel: Erste Hilfe erforderlich & 1\\
\hline
\end{longtable}

\chapter{Chapter 3}
\label{cha:Chapter3}

\section{Section 1}
\label{sec:3Section1}
Willkommen im Portal f"ur Elektronik, Maschinenbau und Mechatronik !
Dieses Portal soll euch beim lernen und diskutieren der einzelnen Studienf"acher behilflich sein, oder im Alltag als Knowledge-Base zur Verf"ugung stehen ! \\
F"ur jedes Studienfach wird in einem "Ubersichtsartikel der Inhalt zusammengefasst und die einzelnen Fachartikel in Beziehung zueinader gestellt. Kommen Formeln in den Fachartikeln vor, werden diese in einer Formelsammlung zu dem jeweiligen Studienfach zusammengef"uhrt. Am Ende soll jedes Studienfach einen "Ubersichtsartikel und wenn m"oglich eine Formelsammlung besitzen.

\subsection{Subcestion 1.1}
\label{subsec:3Subcestion1.1}

\begin{figure}[htb]
\centering
\includegraphics[width=0.2\textwidth]{mkDoc-Logo.png}
\caption{mkDoc}
\label{fig:mkDoc}
\end{figure}


\subsection{Subcestion 1.2}
\label{subsec:3Subcestion 1.2}
Willkommen im Portal f"ur Elektronik, Maschinenbau und Mechatronik !\footnote{\Vgl\Zitat[S.~11]{Sicherheitstechnik}}
Dieses Portal soll euch beim lernen und diskutieren der einzelnen Studienf"acher behilflich sein, oder im Alltag als Knowledge-Base zur Verf"ugung stehen ! \\
F"ur jedes Studienfach wird in einem "Ubersichtsartikel der Inhalt zusammengefasst und die einzelnen Fachartikel in Beziehung zueinader gestellt. Kommen Formeln in den Fachartikeln vor, werden diese in einer Formelsammlung zu dem jeweiligen Studienfach zusammengef"uhrt. Am Ende soll jedes Studienfach einen "Ubersichtsartikel und wenn m"oglich eine Formelsammlung besitzen.

\section{Section 2}
\label{sec:3Section2}

Willkommen im Portal f"ur Elektronik, Maschinenbau und Mechatronik !
Dieses Portal soll euch beim lernen und diskutieren der einzelnen Studienf"acher behilflich sein, oder im Alltag als Knowledge-Base zur Verf"ugung stehen ! \ref{sec:3Section1} \\
F"ur jedes Studienfach wird in einem "Ubersichtsartikel der Inhalt zusammengefasst und die einzelnen Fachartikel in Beziehung zueinader gestellt. Kommen Formeln in den Fachartikeln vor, werden diese in einer Formelsammlung zu dem jeweiligen Studienfach zusammengef"uhrt. Am Ende soll jedes Studienfach einen "Ubersichtsartikel und wenn m"oglich eine Formelsammlung besitzen.

\section{Section 3}
\label{sec:3Section3}

\subsection{Subcestion 3.1}
\label{subsec:3Subcestion3.1}
Willkommen im Portal f"ur Elektronik, Maschinenbau und Mechatronik !
Dieses Portal soll euch beim lernen und diskutieren der einzelnen Studienf"acher behilflich sein, oder im Alltag als Knowledge-Base zur Verf"ugung stehen ! \\
F"ur jedes Studienfach wird in einem "Ubersichtsartikel der Inhalt zusammengefasst und die einzelnen Fachartikel in Beziehung zueinader gestellt. Kommen Formeln in den Fachartikeln vor, werden diese in einer Formelsammlung zu dem jeweiligen Studienfach zusammengef"uhrt. Am Ende soll jedes Studienfach einen "Ubersichtsartikel und wenn m"oglich eine Formelsammlung besitzen. \acs{MTTF}

\subsection{Subcestion 3.2}
\label{subsec:3Subcestion3.2}
Willkommen im Portal f"ur Elektronik, Maschinenbau und Mechatronik !
Dieses Portal soll euch beim lernen und diskutieren der einzelnen Studienf"acher behilflich sein, oder im Alltag als Knowledge-Base zur Verf"ugung stehen ! \\
F"ur jedes Studienfach wird in einem "Ubersichtsartikel der Inhalt zusammengefasst und die einzelnen Fachartikel in Beziehung zueinader gestellt. Kommen Formeln in den Fachartikeln vor, werden diese in einer Formelsammlung zu dem jeweiligen Studienfach zusammengef"uhrt. Am Ende soll jedes Studienfach einen "Ubersichtsartikel und wenn m"oglich eine Formelsammlung besitzen.

\newpage
\subsection{Subcestion 3.3}
\label{subsec:3Subcestion3.3}

\begin{longtable}{|p{9cm}|c|}
\caption{Klassifikation der Schwere (S)}\footnote{\Vgl\Zitat[S.~76]{DINEN62061}}\\
\hline
\label{tab:KlassifikationSchwere}
\textbf{Auswirkungen} & \textbf{Schwere (S)}\\
\hline
\hline
irreversibel: Tod, Verlust eines Auges oder Arms & 4\\
\hline
irreversibel: gebrochene Gliedma\ss en, Verlust (eines) mehrerer Finger(s) & 3\\
\hline
reversibel: Behandlung durch einen Mediziner erforderlich & 2\\
\hline
reversibel: Erste Hilfe erforderlich & 1\\
\hline
\end{longtable}

\chapter{Chapter 4}
\label{cha:Chapter4}

\section{Section 1}
\label{sec:4Section1}
Willkommen im Portal f"ur Elektronik, Maschinenbau und Mechatronik !
Dieses Portal soll euch beim lernen und diskutieren der einzelnen Studienf"acher behilflich sein, oder im Alltag als Knowledge-Base zur Verf"ugung stehen ! \\
F"ur jedes Studienfach wird in einem "Ubersichtsartikel der Inhalt zusammengefasst und die einzelnen Fachartikel in Beziehung zueinader gestellt. Kommen Formeln in den Fachartikeln vor, werden diese in einer Formelsammlung zu dem jeweiligen Studienfach zusammengef"uhrt. Am Ende soll jedes Studienfach einen "Ubersichtsartikel und wenn m"oglich eine Formelsammlung besitzen.

\subsection{Subcestion 1.1}
\label{subsec:4Subcestion1.1}

\begin{figure}[htb]
\centering
\includegraphics[width=0.2\textwidth]{mkDoc-Logo.png}
\caption{mkDoc}
\label{fig:mkDoc}
\end{figure}


\subsection{Subcestion 1.2}
\label{subsec:4Subcestion 1.2}
Willkommen im Portal f"ur Elektronik, Maschinenbau und Mechatronik !\footnote{\Vgl\Zitat[S.~11]{Sicherheitstechnik}}
Dieses Portal soll euch beim lernen und diskutieren der einzelnen Studienf"acher behilflich sein, oder im Alltag als Knowledge-Base zur Verf"ugung stehen ! \\
F"ur jedes Studienfach wird in einem "Ubersichtsartikel der Inhalt zusammengefasst und die einzelnen Fachartikel in Beziehung zueinader gestellt. Kommen Formeln in den Fachartikeln vor, werden diese in einer Formelsammlung zu dem jeweiligen Studienfach zusammengef"uhrt. Am Ende soll jedes Studienfach einen "Ubersichtsartikel und wenn m"oglich eine Formelsammlung besitzen.

\section{Section 2}
\label{sec:4Section2}

Willkommen im Portal f"ur Elektronik, Maschinenbau und Mechatronik !
Dieses Portal soll euch beim lernen und diskutieren der einzelnen Studienf"acher behilflich sein, oder im Alltag als Knowledge-Base zur Verf"ugung stehen ! \ref{sec:4Section1} \\
F"ur jedes Studienfach wird in einem "Ubersichtsartikel der Inhalt zusammengefasst und die einzelnen Fachartikel in Beziehung zueinader gestellt. Kommen Formeln in den Fachartikeln vor, werden diese in einer Formelsammlung zu dem jeweiligen Studienfach zusammengef"uhrt. Am Ende soll jedes Studienfach einen "Ubersichtsartikel und wenn m"oglich eine Formelsammlung besitzen.

\section{Section 3}
\label{sec:4Section3}

\subsection{Subcestion 3.1}
\label{subsec:4Subcestion3.1}
Willkommen im Portal f"ur Elektronik, Maschinenbau und Mechatronik !
Dieses Portal soll euch beim lernen und diskutieren der einzelnen Studienf"acher behilflich sein, oder im Alltag als Knowledge-Base zur Verf"ugung stehen ! \\
F"ur jedes Studienfach wird in einem "Ubersichtsartikel der Inhalt zusammengefasst und die einzelnen Fachartikel in Beziehung zueinader gestellt. Kommen Formeln in den Fachartikeln vor, werden diese in einer Formelsammlung zu dem jeweiligen Studienfach zusammengef"uhrt. Am Ende soll jedes Studienfach einen "Ubersichtsartikel und wenn m"oglich eine Formelsammlung besitzen. \acs{MTTF}

\subsection{Subcestion 3.2}
\label{subsec:4Subcestion3.2}
Willkommen im Portal f"ur Elektronik, Maschinenbau und Mechatronik !
Dieses Portal soll euch beim lernen und diskutieren der einzelnen Studienf"acher behilflich sein, oder im Alltag als Knowledge-Base zur Verf"ugung stehen ! \\
F"ur jedes Studienfach wird in einem "Ubersichtsartikel der Inhalt zusammengefasst und die einzelnen Fachartikel in Beziehung zueinader gestellt. Kommen Formeln in den Fachartikeln vor, werden diese in einer Formelsammlung zu dem jeweiligen Studienfach zusammengef"uhrt. Am Ende soll jedes Studienfach einen "Ubersichtsartikel und wenn m"oglich eine Formelsammlung besitzen.

\newpage
\subsection{Subcestion 3.3}
\label{subsec:4Subcestion3.3}

\begin{longtable}{|p{9cm}|c|}
\caption{Klassifikation der Schwere (S)}\footnote{\Vgl\Zitat[S.~76]{DINEN62061}}\\
\hline
\label{tab:KlassifikationSchwere}
\textbf{Auswirkungen} & \textbf{Schwere (S)}\\
\hline
\hline
irreversibel: Tod, Verlust eines Auges oder Arms & 4\\
\hline
irreversibel: gebrochene Gliedma\ss en, Verlust (eines) mehrerer Finger(s) & 3\\
\hline
reversibel: Behandlung durch einen Mediziner erforderlich & 2\\
\hline
reversibel: Erste Hilfe erforderlich & 1\\
\hline
\end{longtable}

\chapter{Chapter 5}
\label{cha:Chapter5}

\section{Section 1}
\label{sec:5Section1}
Willkommen im Portal f"ur Elektronik, Maschinenbau und Mechatronik !
Dieses Portal soll euch beim lernen und diskutieren der einzelnen Studienf"acher behilflich sein, oder im Alltag als Knowledge-Base zur Verf"ugung stehen ! \\
F"ur jedes Studienfach wird in einem "Ubersichtsartikel der Inhalt zusammengefasst und die einzelnen Fachartikel in Beziehung zueinader gestellt. Kommen Formeln in den Fachartikeln vor, werden diese in einer Formelsammlung zu dem jeweiligen Studienfach zusammengef"uhrt. Am Ende soll jedes Studienfach einen "Ubersichtsartikel und wenn m"oglich eine Formelsammlung besitzen.

\subsection{Subcestion 1.1}
\label{subsec:5Subcestion1.1}

\begin{figure}[htb]
\centering
\includegraphics[width=0.2\textwidth]{mkDoc-Logo.png}
\caption{mkDoc}
\label{fig:mkDoc}
\end{figure}


\subsection{Subcestion 1.2}
\label{subsec:5Subcestion 1.2}
Willkommen im Portal f"ur Elektronik, Maschinenbau und Mechatronik !\footnote{\Vgl\Zitat[S.~11]{Sicherheitstechnik}}
Dieses Portal soll euch beim lernen und diskutieren der einzelnen Studienf"acher behilflich sein, oder im Alltag als Knowledge-Base zur Verf"ugung stehen ! \\
F"ur jedes Studienfach wird in einem "Ubersichtsartikel der Inhalt zusammengefasst und die einzelnen Fachartikel in Beziehung zueinader gestellt. Kommen Formeln in den Fachartikeln vor, werden diese in einer Formelsammlung zu dem jeweiligen Studienfach zusammengef"uhrt. Am Ende soll jedes Studienfach einen "Ubersichtsartikel und wenn m"oglich eine Formelsammlung besitzen.

\section{Section 2}
\label{sec:5Section2}

Willkommen im Portal f"ur Elektronik, Maschinenbau und Mechatronik !
Dieses Portal soll euch beim lernen und diskutieren der einzelnen Studienf"acher behilflich sein, oder im Alltag als Knowledge-Base zur Verf"ugung stehen ! \ref{sec:5Section1} \\
F"ur jedes Studienfach wird in einem "Ubersichtsartikel der Inhalt zusammengefasst und die einzelnen Fachartikel in Beziehung zueinader gestellt. Kommen Formeln in den Fachartikeln vor, werden diese in einer Formelsammlung zu dem jeweiligen Studienfach zusammengef"uhrt. Am Ende soll jedes Studienfach einen "Ubersichtsartikel und wenn m"oglich eine Formelsammlung besitzen.

\section{Section 3}
\label{sec:5Section3}

\subsection{Subcestion 3.1}
\label{subsec:4Subcestion3.1}
Willkommen im Portal f"ur Elektronik, Maschinenbau und Mechatronik !
Dieses Portal soll euch beim lernen und diskutieren der einzelnen Studienf"acher behilflich sein, oder im Alltag als Knowledge-Base zur Verf"ugung stehen ! \\
F"ur jedes Studienfach wird in einem "Ubersichtsartikel der Inhalt zusammengefasst und die einzelnen Fachartikel in Beziehung zueinader gestellt. Kommen Formeln in den Fachartikeln vor, werden diese in einer Formelsammlung zu dem jeweiligen Studienfach zusammengef"uhrt. Am Ende soll jedes Studienfach einen "Ubersichtsartikel und wenn m"oglich eine Formelsammlung besitzen. \acs{MTTF}

\subsection{Subcestion 3.2}
\label{subsec:5Subcestion3.2}
Willkommen im Portal f"ur Elektronik, Maschinenbau und Mechatronik !
Dieses Portal soll euch beim lernen und diskutieren der einzelnen Studienf"acher behilflich sein, oder im Alltag als Knowledge-Base zur Verf"ugung stehen ! \\
F"ur jedes Studienfach wird in einem "Ubersichtsartikel der Inhalt zusammengefasst und die einzelnen Fachartikel in Beziehung zueinader gestellt. Kommen Formeln in den Fachartikeln vor, werden diese in einer Formelsammlung zu dem jeweiligen Studienfach zusammengef"uhrt. Am Ende soll jedes Studienfach einen "Ubersichtsartikel und wenn m"oglich eine Formelsammlung besitzen.

\newpage
\subsection{Subcestion 3.3}
\label{subsec:5Subcestion3.3}

\begin{longtable}{|p{9cm}|c|}
\caption{Klassifikation der Schwere (S)}\footnote{\Vgl\Zitat[S.~76]{DINEN62061}}\\
\hline
\label{tab:KlassifikationSchwere}
\textbf{Auswirkungen} & \textbf{Schwere (S)}\\
\hline
\hline
irreversibel: Tod, Verlust eines Auges oder Arms & 4\\
\hline
irreversibel: gebrochene Gliedma\ss en, Verlust (eines) mehrerer Finger(s) & 3\\
\hline
reversibel: Behandlung durch einen Mediziner erforderlich & 2\\
\hline
reversibel: Erste Hilfe erforderlich & 1\\
\hline
\end{longtable}

%\include{Inhalt/ZitateReferenzen}
%\chapter{Bilder und Listings}

\section{pasCAL}
\label{sec:pasCAL}

\section{SISTEMA}
\label{sec:SISTEMA}


\chapter{Kritische Bewertung}
\label{cha:Fazit}

\section{Section 1}
\label{sec:6Section1}
Willkommen im Portal f"ur Elektronik, Maschinenbau und Mechatronik !
Dieses Portal soll euch beim lernen und diskutieren der einzelnen Studienf"acher behilflich sein, oder im Alltag als Knowledge-Base zur Verf"ugung stehen ! \\
F"ur jedes Studienfach wird in einem "Ubersichtsartikel der Inhalt zusammengefasst und die einzelnen Fachartikel in Beziehung zueinader gestellt. Kommen Formeln in den Fachartikeln vor, werden diese in einer Formelsammlung zu dem jeweiligen Studienfach zusammengef"uhrt. Am Ende soll jedes Studienfach einen "Ubersichtsartikel und wenn m"oglich eine Formelsammlung besitzen.

\subsection{Subcestion 1.1}
\label{subsec:6Subcestion1.1}

\begin{figure}[htb]
\centering
\includegraphics[width=0.2\textwidth]{mkDoc-Logo.png}
\caption{mkDoc}
\label{fig:mkDoc}
\end{figure}


\subsection{Subcestion 1.2}
\label{subsec:6Subcestion 1.2}
Willkommen im Portal f"ur Elektronik, Maschinenbau und Mechatronik !\footnote{\Vgl\Zitat[S.~11]{Sicherheitstechnik}}
Dieses Portal soll euch beim lernen und diskutieren der einzelnen Studienf"acher behilflich sein, oder im Alltag als Knowledge-Base zur Verf"ugung stehen ! \\
F"ur jedes Studienfach wird in einem "Ubersichtsartikel der Inhalt zusammengefasst und die einzelnen Fachartikel in Beziehung zueinader gestellt. Kommen Formeln in den Fachartikeln vor, werden diese in einer Formelsammlung zu dem jeweiligen Studienfach zusammengef"uhrt. Am Ende soll jedes Studienfach einen "Ubersichtsartikel und wenn m"oglich eine Formelsammlung besitzen.



\chapter{Zusammenfassung}
\label{cha:Zusammenfassung}

\section{Section 1}
\label{sec:7Section1}
Willkommen im Portal f"ur Elektronik, Maschinenbau und Mechatronik !
Dieses Portal soll euch beim lernen und diskutieren der einzelnen Studienf"acher behilflich sein, oder im Alltag als Knowledge-Base zur Verf"ugung stehen ! \\
F"ur jedes Studienfach wird in einem "Ubersichtsartikel der Inhalt zusammengefasst und die einzelnen Fachartikel in Beziehung zueinader gestellt. Kommen Formeln in den Fachartikeln vor, werden diese in einer Formelsammlung zu dem jeweiligen Studienfach zusammengef"uhrt. Am Ende soll jedes Studienfach einen "Ubersichtsartikel und wenn m"oglich eine Formelsammlung besitzen.

\subsection{Subcestion 1.1}
\label{subsec:7Subcestion1.1}

\begin{figure}[htb]
\centering
\includegraphics[width=0.2\textwidth]{mkDoc-Logo.png}
\caption{mkDoc}
\label{fig:mkDoc}
\end{figure}


\subsection{Subcestion 1.2}
\label{subsec:7Subcestion 1.2}
Willkommen im Portal f"ur Elektronik, Maschinenbau und Mechatronik !\footnote{\Vgl\Zitat[S.~11]{Sicherheitstechnik}}
Dieses Portal soll euch beim lernen und diskutieren der einzelnen Studienf"acher behilflich sein, oder im Alltag als Knowledge-Base zur Verf"ugung stehen ! \\
F"ur jedes Studienfach wird in einem "Ubersichtsartikel der Inhalt zusammengefasst und die einzelnen Fachartikel in Beziehung zueinader gestellt. Kommen Formeln in den Fachartikeln vor, werden diese in einer Formelsammlung zu dem jeweiligen Studienfach zusammengef"uhrt. Am Ende soll jedes Studienfach einen "Ubersichtsartikel und wenn m"oglich eine Formelsammlung besitzen.

\chapter{Vom Autor verwendete Software}
\label{cha:Werkzeuge}

\begin{itemize}
\itemd{Inkscape}{Version: 0.47,
\newline Website: \url{http://www.inkscape.org/download/?lang=de},
\newline Stand: 09.08.2009.}
\itemd{LTi SafePLC}{Version: 1.25,
\newline Website: \url{http://www.lust-antriebstechnik.de/},
\newline Stand: 09.11.2009.}
\itemd{Pilz PAScal}{Version: v1.3.2 Build 004,
\newline Website: \url{http://www.pilz.com/downloads/restricted/pascal_1.3.2.zip},
\newline Stand: 14.10.2009.}
\itemd{SISTEMA}{Version: 1.0.5,\newline Website: \url{http://www.dguv.de/bgia/de/pra/softwa/sistema/index.jsp},
\newline Stand: 14.10.2009.}
\itemd{TeXShop \LaTeX }{Version: 2.26 Release 3/17/2009,
\newline Website: \url{http://www.uoregon.edu/~koch/texshop/index.html},
\newline Stand: 12.09.2009.}
\end{itemize}



% Literaturverzeichnis ---------------------------------------------------------
%   Das Literaturverzeichnis wird aus der BibTeX-Datenbank "Bibliographie.bib"
%   erstellt.
% ------------------------------------------------------------------------------
\bibliography{Bibliographie} % Aufruf: bibtex Diplomarbeit
\bibliographystyle{natdin} % DIN-Stil des Literaturverzeichnisses

% Index ------------------------------------------------------------------------
%   Zum Erstellen eines Index, die folgende Zeile auskommentieren.
% ------------------------------------------------------------------------------
%\printindex

\addchap{Eidesstattliche Erkl"arung}
Ich, \autor, Matrikel-Nr.\ \matrikelnr, versichere hiermit, dass ich meine Diplomarbeit mit dem Thema
\begin{quote}
\textit{\titel} \textit{\untertitel}
\end{quote}
selbstst"andig verfasst und keine anderen als die angegebenen Quellen und Hilfsmittel benutzt habe, wobei ich alle w�rtlichen und sinngem"a\ss en Zitate als solche gekennzeichnet habe. Die Arbeit wurde bisher keiner anderen Pr"ufungsbeh"orde vorgelegt und auch nicht ver"offentlicht.

Mir ist bekannt, dass ich meine Diplomarbeit zusammen mit dieser Erkl"arung fristgem"a\ss \ nach Vergabe des Themas in dreifacher Ausfertigung und gebunden im Pr"ufungsamt der Wilhelm-B\"uchner-Hochschule abzugeben oder sp"atestens mit dem Poststempel des Tages, an dem die Frist abl"auft, zu senden habe.\\[6ex]

\ort, den 
\newline

\rule[-0.2cm]{5cm}{0.5pt}

\textsc{\autor} 
 % Selbst�ndigkeitserkl�rung

% Anhang -----------------------------------------------------------------------
%   Die Inhalte des Anhangs werden analog zu den Kapiteln inkludiert.
%   Dies geschieht in der Datei "Anhang.tex".
% ------------------------------------------------------------------------------
\begin{appendix}
    \clearpage
    \pagenumbering{roman}
    \chapter{Anhang}
    \label{sec:Anhang}
    % Rand der Aufz�hlungen in Tabellen anpassen
    \setdefaultleftmargin{1em}{}{}{}{}{}
    \section{DVD}
\label{sec:DVD}

\begin{description}
\item[Diplomarbeit]\begin{it}PDF-Datei und Anh"ange \end{it}
\begin{itemize}
\item{Diplomarbeit} 
\item{Dokument1}
\item{Dokument2} 
\item{Dokument3}
\item{Dokument4}
\item{Dokument5}
\item{Dokument6}
\item{Dokument7}
\end{itemize}
\item[Sonstiges]\begin{it}Zus"atzliche Informationonen zur Diplomarbeit \end{it}
\begin{itemize}
\item{Dokument1}
\item{Dokument2} 
\item{Dokument3}
\end{itemize}
\end{description}



\section{mkDoc2}
\label{sec:mkDoc2}
\vspace{4mm}
\includegraphics[width=0.3\textwidth]{mkDoc-Logo.png}

\section{mkDoc3}
\label{sec:mkDoc3}
\vspace{4mm}
\includegraphics[width=0.3\textwidth]{mkDoc-Logo.png}


\end{appendix}


\end{document}
